% Source: Weinberg GR, physical PDF pp. 28--29
%         (printed pp. 3--4).
% Coverage: Chapter 1 opening epigraph and introductory prose.
% Figures/tables/footnotes: no figures or tables; reference marker 1.
% Status: source-reviewed and compile-clean.
% Uncertainties: none.

\begin{flushright}
\begin{minipage}{0.43\textwidth}
\small
``But the tale of history forms a very strong bulwark against the stream of
time, and to some extent checks its irresistible flow, and, of all things done
in it, as many as history has taken over, it secures and binds together, and
does not allow them to slip away into the abyss of oblivion.''

\medskip
\emph{Anna Comnena, The Alexiad}
\end{minipage}
\end{flushright}

Physics is not a finished logical system. Rather, at any moment it spans a
great confusion of ideas, some that survive like folk epics from the heroic
periods of the past, and others that arise like utopian novels from our dim
premonitions of a future grand synthesis. The author of a book on physics can
impose order on this confusion by organizing his material in either of two
ways: by recapitulating its history, or by following his own best guess as to
the ultimate logical structure of physical law. Both methods are valuable; the
great thing is not to confuse physics with history, or history with physics.

This book sets out the theory of gravitation according to what I think is its
inner logic as a branch of physics, and not according to its historical
development. It is certainly a historical fact that when Albert Einstein was
working out general relativity, there was at hand a preexisting mathematical
formalism, that of Riemannian geometry, that he could and did take over whole.
However, this historical fact does not mean that the essence of general
relativity necessarily consists in the application of Riemannian geometry to
physical space and time. In my view, it is much more useful to regard general
relativity above all as a theory of \emph{gravitation}, whose connection with
geometry arises from the peculiar empirical properties of gravitation,
properties summarized by Einstein's Principle of the Equivalence of
Gravitation and Inertia. For this reason, I have tried throughout this book to
delay the introduction of geometrical objects, such as the metric, the affine
connection, and the curvature, until the use of these objects could be
motivated by considerations of physics. The order of chapters here thus bears
very little resemblance to the order of history.

Nevertheless, because we must not allow the history of physics ``to slip away
into the abyss of oblivion,'' this first chapter presents a brief backward look
at three great antecedents to general relativity---non-Euclidean geometry, the
Newtonian theory of gravitation, and the principle of relativity. Their history
is traced up to 1916, the year in which they were brought together by Einstein
in the General Theory of
Relativity.\hyperref[ch1-ref-1]{\textsuperscript{1}}
